\subsection{Tensor products and homotopy}

PROPOSITION 3. --- Let C, C \(_{1}\) be two complexes of right A-modules, C', C \(_{1}'\) two complexes of left A-modules, and u : C → C \(_{1}\) , v : C → C \(_{1}\) , u' : C' → C \(_{1}'\) , v' : C' → C \(_{1}'\) be morphisms of complexes.

\begin{enumerate}
\def\labelenumi{\alph{enumi})}
\item
  If \(u\) and \(u'\) are homotopic to \(v\) and \(v'\) respectively, then the two morphisms \(u \otimes u'\) and \(v \otimes v'\) from \(\mathbf{C} \otimes_{\mathbf{A}} \mathbf{C}'\) to \(\mathbf{C}_1 \otimes_{\mathbf{A}} \mathbf{C}_1'\) are homotopic.
\item
  If u and \(u'\) are homotopisms, \(u \otimes u'\) is a homotopism.
\item
  If C or \(\mathbf{C}'\) is homotopic to zero, \(\mathbf{C} \otimes_{\mathbf{A}} \mathbf{C}'\) is homotopic to zero.
\end{enumerate}

Let us denote by the same letter d the differentials of the complexes C, \(C_{1}\) , \(C'\) , \(C_{1}'\) and by D the differentials of the complexes \(C \otimes_{A} C'\) and \(C_{1} \otimes_{A} C_{1}'\) .

If \(u\) (resp. \(u'\) ) is homotopic to \(v\) (resp. \(v'\) ), there exists a graded homomorphism of degree 1 \(s: \mathbf{C} \to \mathbf{C}_1\) (resp. \(s': \mathbf{C}' \to \mathbf{C}_1'\) ) such that

\[
u - v = d s + s d \qquad (\mathrm{resp.} u ^ {\prime} - v ^ {\prime} = d s ^ {\prime} + s ^ {\prime} d).\tag{2}
\]

Let \(\mathbf{S}:\mathbf{C}\otimes_{\mathbf{A}}\mathbf{C}^{\prime}\to \mathbf{C}_{1}\otimes_{\mathbf{A}}\mathbf{C}_{1}^{\prime}\) be the unique graded homomorphism of degree 1 such that, for \(x\in \mathbf{C}_p\) , \(y\in \mathbf{C}_q^\prime\) , \(p,q\in \mathbf{Z}\) , one has

\[
\mathbf {S} (x \otimes y) = s (x) \otimes u ^ {\prime} (y) + (- 1) ^ {p} v (x) \otimes s ^ {\prime} (y).\tag{3}
\]

We then have, with the preceding notation:

\[
\begin{array}{r l} (\mathrm{DS} + \mathrm{SD}) (x \otimes y) & = \mathrm{D} (s x \otimes u ^ {\prime} y) + (- 1) ^ {p} \mathrm{D} (v x \otimes s ^ {\prime} y) + \mathrm{S} (d x \otimes y) + \\ & + (- 1) ^ {p} \mathrm{S} (x \otimes d y) = \\ & = d s x \otimes u ^ {\prime} y + (- 1) ^ {p + 1} s x \otimes d u ^ {\prime} y + (- 1) ^ {p} d v x \otimes s ^ {\prime} y + v x \otimes d s ^ {\prime} y + \\ & + s d x \otimes u ^ {\prime} y + (- 1) ^ {p - 1} v d x \otimes s ^ {\prime} y + (- 1) ^ {p} s x \otimes u ^ {\prime} d y + v x \otimes s ^ {\prime} d y \\ & = (d s + s d) (x) \otimes u ^ {\prime} y + v x \otimes (d s ^ {\prime} + s ^ {\prime} d) (y) \\ & = (u x - v x) \otimes u ^ {\prime} y + v x \otimes (u ^ {\prime} y - v ^ {\prime} y) = u x \otimes u ^ {\prime} y - v x \otimes v ^ {\prime} y. \end{array}
\]

This gives DS + SD = u ⊗ u' - v ⊗ v', whence a).

Let us prove b). If u and \(u'\) are homotopisms, there exist homomorphisms of complexes \(\alpha: C_{1} \to C\) and \(\alpha': C_{1}' \to C'\) such that \(u \circ \alpha\) , \(\alpha \circ u\) , \(u' \circ \alpha'\) , \(\alpha' \circ u'\) are homotopic respectively to \(Id_{C_{1}}\) , \(Id_{C}\) , \(Id_{C_{1}}\) and \(Id_{C'}\) . Then \((u \otimes u') \circ (\alpha \otimes \alpha')\) , which is equal to \((u \circ \alpha) \otimes (u' \circ \alpha')\) , is homotopic by a) to \(Id_{C_{1}} \otimes Id_{C_{1}} = Id_{C_{1} \otimes C_{1}}\) , whereas \((\alpha \otimes \alpha') \circ (u \otimes u')\) is homotopic to \(Id_{C \otimes C'}\) , whence b). Finally, c) follows from b) applied to the case where \(C_{1}\) or \(C_{1}'\) is zero.

COROLLARY 1. --- Let C' be a split complex of left A-modules such that H(C') is flat. For every complex C of right A-modules, the canonical map

\[
\gamma (\mathrm{C}, \mathrm{C} ^ {\prime}): \mathrm{H} (\mathrm{C}) \otimes_ {\mathrm{A}} \mathrm{H} (\mathrm{C} ^ {\prime}) \rightarrow \mathrm{H} (\mathrm{C} \otimes_ {\mathrm{A}} \mathrm{C} ^ {\prime})
\]

is bijective.

According to X, p.~35, def. 6, there exists a homotopism \(u': C' \to H(C')\) . By prop. 3, \(1_{C} \otimes u': C \otimes_{A} C' \to C \otimes_{A} H(C')\) is a homotopism; since

\[
\mathrm{H} \left(1 _ {\mathrm{C}} \otimes u ^ {\prime}\right) \circ \gamma (\mathrm{C}, \mathrm{C} ^ {\prime}) = \gamma (\mathrm{C}, \mathrm{H} (\mathrm{C} ^ {\prime})) \circ \left(1 _ {\mathrm{C}} \otimes \mathrm{H} (u ^ {\prime})\right),
\]

and \(\mathbf{H}(1_{\mathbf{C}} \otimes u')\) and \(\mathbf{H}(u')\) are bijective, it suffices to prove that \(\gamma(\mathbf{C}, \mathbf{H}(\mathbf{C}'))\) is bijective, and we are reduced to the case where \(\mathbf{C}'\) is flat and has zero differential. In this case the canonical exact sequences

\begin{enumerate}
\def\labelenumi{(\Roman{enumi})}
\tightlist
\item
\end{enumerate}

\[
0 \to \mathbf {Z} (\mathbf {C}) \xrightarrow {i} \mathbf {C} \xrightarrow {\delta} \mathbf {B} (\mathbf {C}) \to 0\tag{II}
\]

\[
0 \rightarrow \mathrm{B} (\mathrm{C}) \xrightarrow {j} \mathrm{Z} (\mathrm{C}) \xrightarrow {\pi} \mathrm{H} (\mathrm{C}) \rightarrow 0
\]

give exact sequences:

\[
\begin{array}{l} 0 \longrightarrow Z (C) \otimes_ {A} C ^ {\prime} \xrightarrow {i \otimes 1} C \otimes_ {A} C ^ {\prime} \quad \xrightarrow {\delta \otimes 1} B (C) \otimes_ {A} C ^ {\prime} \longrightarrow 0 \\ 0 \longrightarrow B (C) \otimes_ {A} C ^ {\prime} \xrightarrow {j \otimes 1} Z (C) \otimes_ {A} C ^ {\prime} \xrightarrow {\pi \otimes 1} H (C) \otimes_ {A} C ^ {\prime} \longrightarrow 0. \end{array}
\]

Since \(d = i \circ j \circ \delta\) , we have \(\mathbf{D} = d \otimes 1_{\mathbf{C}'} = (i \otimes 1) \circ (j \otimes 1) \circ (\delta \otimes 1)\), which shows that the canonical maps \(\mathbf{Z}(\mathbf{C}) \otimes_{\mathbf{A}} \mathbf{C}' \to \mathbf{Z}(\mathbf{C} \otimes_{\mathbf{A}} \mathbf{C}')\) and \(\mathbf{B}(\mathbf{C}) \otimes_{\mathbf{A}} \mathbf{C}' \to \mathbf{B}(\mathbf{C} \otimes_{\mathbf{A}} \mathbf{C}')\) are bijective, hence also \(\gamma(\mathbf{C}, \mathbf{C}')\) by passing to quotients.

COROLLARY 2. --- Let N be a flat left A-module. For every complex C of right A-modules, the canonical homomorphisms

\[
\gamma_ {n} (\mathbf {C}, \mathbf {N}): \mathrm{H} _ {n} (\mathbf {C}) \otimes_ {\mathbf {A}} \mathbf {N} \rightarrow \mathrm{H} _ {n} (\mathbf {C} \otimes_ {\mathbf {A}} \mathbf {N})
\]

are bijective.

COROLLARY 3. --- Let C' be a complex of left A-modules such that B(C') and H(C') are projective. For every complex C of right A-modules, the map γ(C, C') is bijective.

Indeed, C' is split (X, p.~35, example 4) and H(C') is projective; one can therefore apply Corollary 1.

Remarks. --- 1) Using the commutation isomorphisms, one deduces from Corollaries 1, 2, and 3 the analogous statements obtained by interchanging the roles of the two arguments of the tensor products.

\begin{enumerate}
\def\labelenumi{\arabic{enumi})}
\setcounter{enumi}{1}
\tightlist
\item
  We shall see below (X, p.~79, cor. 4) that the conclusion of cor. 1 is also true when one assumes that C' and H(C') are flat and C' is bounded on the right.
\end{enumerate}

